% Draft Related Work section. All keys are defined in references_verified.bib.
\section{Related Work}
\label{sec:related}

\textbf{Perishable inventory control.} Fixed-shelf-life inventory has a long theoretical tradition \cite{nahmias1982perishable}. Optimal policies are hard to compute because the state must record the age profile of the stock, which has motivated near-myopic and heuristic policies \cite{nandakumar1993near}. Broekmeulen and van Donselaar \cite{broekmeulen2009heuristic} give a heuristic for the retail case with batch ordering, positive lead times and time-varying demand, and Haijema and Minner \cite{haijema2019improved} quantify the value of using stock-age information when ordering. Reviews of deteriorating-inventory models cover the wider literature \cite{bakker2012review, janssen2016literature}. We do not propose a new policy. We fix a simple order-up-to rule with first-in-first-out issuing and use it as a common yardstick for the forecasts.

\textbf{Calendar, event and promotion effects in retail forecasting.} Holidays, promotions and other special days are among the main drivers of retail demand, and their treatment is a recurring topic in reviews of retail forecasting \cite{fildes2022retail}. Promotional forecasting has been studied with regression-type models \cite{trapero2015identification}, including for perishable items, where promotional uplift interacts with waste \cite{vandonselaar2016analysis}. Huber and Stuckenschmidt \cite{huber2020daily} compare machine-learning forecasters for daily retail demand with a focus on calendric special days. Public benchmarks supply event calendars as features: the Favorita data include a holiday calendar \cite{favorita2018kaggle}, and the M5 data include calendar events and SNAP days \cite{makridakis2022m5background}. In M5, methods built on gradient-boosted trees \cite{friedman2001gbm, ke2017lightgbm} performed strongly \cite{makridakis2022m5}. External signals such as weather have also been shown to carry operational value \cite{steinker2017value}. In these studies the event inputs are usually treated as known and correct; how errors in them affect results is rarely examined.

\textbf{Forecast quality judged by inventory outcomes.} Several authors argue that forecasts for replenishment should be judged by the inventory performance they produce, not only by accuracy. Gardner \cite{gardner1990evaluating} evaluates forecasting methods through inventory trade-off curves, and Kourentzes et al. \cite{kourentzes2020optimising} show that standard accuracy measures can be poorly aligned with inventory outcomes and propose optimising forecasting models for inventory planning. A related line connects forecasting and decisions directly. Quantile regression \cite{koenker1978regression} gives the critical-fractile quantile that a newsvendor order requires, and feature-based or data-driven newsvendor models learn order quantities from covariates \cite{ban2019big, huber2019data, bertsimas2020predictive}. Our forecasters follow this pattern: a gradient-boosted quantile model sets the order-up-to level, and the arms are compared by simulated cost.

\textbf{Imperfect inputs and train--serve mismatch.} When the distribution of inputs at deployment differs from that seen in training, a model fitted on one can perform poorly on the other; this is the general problem of dataset shift \cite{quinonero2009dataset}. A clean, after-the-fact event calendar used for training but not available at ordering time is closely related to data leakage, in which a model is built with information that will not exist when predictions are made \cite{kaufman2012leakage}. Learning using privileged information \cite{vapnik2009new} studies the opposite design choice, deliberately exploiting information that is present only during training. Two remedies for noisy inputs are also relevant. Adding noise to inputs during training acts as a regulariser \cite{bishop1995training}, and data augmentation for time series is an active area \cite{wen2021time}. Our arm D, trained on the clean calendar with simulated signal errors injected at estimated rates, is an instance of this idea applied to event features.

\textbf{Positioning.} This paper combines these threads in one setting. Rather than asking whether event features improve accuracy, we ask how the quality of the event signal available at ordering time changes the cost of a perishable replenishment policy. We compare training on the clean calendar, on the signals as they would have been available, and on the clean calendar with injected noise, and we vary the error rates used in training and in serving separately. We also vary event strength independently of signal quality, so that results are reported as a map of regimes rather than a single effect size. To our knowledge, the effect of imperfect event signals on perishable inventory cost, and the choice of training signal under such errors, has received little attention in the literature reviewed above.
